\section{A uniform local index}
\label{sec:index}

The small sections constructed above vanish to a definite weighted order
at the marked points.  The essential issue is uniformity: the marked
points can have arbitrarily large degrees, and the coefficients of the
homomorphisms tend to infinity.  We prove the local assertion with these
two features explicit. The weighted Taylor-coefficient and Cauchy
argument follows the method of Faltings \cite[p.~560]{Faltings1991};
the estimates below supply the required uniformity for our varying
varieties and homomorphisms.

For a smooth point $y=(y_\ell)_\ell$ of a product
$Z=\prod_\ell Z_\ell$, choose regular parameters
$T_\ell=(T_{\ell,1},\ldots,T_{\ell,d_\ell})$ on each factor and a
local frame $e$ of a line bundle $E$.  If a nonzero section of $E^n$
has expansion $s=e^n\sum_\alpha a_\alpha T^\alpha$, define
\begin{equation}\label{eq:weighted-index}
 \operatorname{ind}_{n,b}(s;y)
 =\min_{a_\alpha\ne0}
       \sum_\ell\frac{|\alpha_\ell|}{n b_\ell^2}.
\end{equation}
A zero-dimensional factor contributes no parameters.  All parameters
in a given factor have the same weight.  Factorwise changes of regular
parameters preserve the corresponding filtration of the completed
local ring; multiplication by a unit also preserves its first nonzero
graded piece.  Thus the index is independent of these choices.

\begin{proposition}\label{prop:index}
Fix projective embeddings of abelian varieties $A_\ell$ by symmetric
very ample bundles $L_\ell$, a symmetric very ample canonically metrized
bundle $\bar L$ on $\cA=\prod_\ell A_\ell$, and integral bases of the
homomorphism groups between the fixed varieties.  Let
$F_i:\cA\to\cA$ be homomorphisms.  Suppose that:
\begin{enumerate}
\item $y_{i\ell}$ is a smooth point of an irreducible
$Z_{i\ell}\subset A_\ell$, its projective height is $O(H_{i\ell})$,
and its conormal space is spanned by equations of bounded degree and
coefficient height $O(H_{i\ell})$, with $H_{i\ell}\ge1$;
\item $M_i\longrightarrow\infty$, the positive integers $b_{i\ell}$
satisfy $b_{i\ell}\ll M_i$ and
$b_{i\ell}^2H_{i\ell}\asymp M_i^2$, and the coefficients of
$F_i|_{A_\ell}$ in the fixed bases are $O(b_{i\ell})$;
\item for $\bar E_i=F_i^*\bar L$ one has
$h_{\bar E_i}(y_i)=o(M_i^2)$, and a nonzero section $s_i$ of
$E_i^{n_i}$ on a resolution of $Z_i=\prod_\ell Z_{i\ell}$ satisfies
\begin{equation}\label{eq:index-small-section}
 \|s_i\|_{v,\sup}\le
 \begin{cases}
 e^{-c n_iM_i^2},&v\mid\infty,\\
 1,&v\nmid\infty,
 \end{cases}
\end{equation}
for a fixed $c>0$, over a number field containing its data.
\end{enumerate}
Assume that the resolution is an isomorphism over the smooth locus.
Then, for some $\eta>0$ independent of $i$,
$\operatorname{ind}_{n_i,b_i}(s_i;y_i)\ge\eta$ for all sufficiently
large $i$.  No bound on the degree of the field of $y_i$ is required.
\end{proposition}

\begin{proof}
For a first nonzero Taylor coefficient, Cauchy's inequality and the
product formula will compare the cost of differentiation,
\(\sum_\ell|\alpha_\ell|H_{i\ell}\), with the small-section gain
\(n_iM_i^2\).  The relation \(b_{i\ell}^2H_{i\ell}\asymp M_i^2\)
will turn this comparison into the asserted weighted index.  We first
choose the algebraic parameters and coefficient to which this argument
will apply, then bound the analytic radii and metric variation.

\emph{An algebraic coefficient and its local estimates.}
Enlarge the field of the section to a number field $E$ containing
$y_i$ and the equations to be used below.  At a place $v$ of $E$ write
$\nu_v=[E_v:\mathbb Q_{v|\mathbb Q}]/[E:\mathbb Q]$, with the usual
absolute values.  All heights and sums over places below use these
weights.  The constants do not depend on $[E:\mathbb Q]$.

Choose an affine projective-coordinate chart at $y_{i\ell}$.
From the equations in the first hypothesis choose a maximal invertible
Jacobian minor.  Center the coordinates at $y_{i\ell}$, write
$U_\ell$ for the coordinates in that minor and $T_\ell$ for the
remaining coordinates, and multiply the equations by the inverse
minor matrix.  They become
\begin{equation}\label{eq:index-implicit-system}
 U_\ell+\beta_\ell T_\ell+Q_\ell(T_\ell,U_\ell)=0,
\end{equation}
where $Q_\ell$ has no terms of degree less than two.  All degrees and
the number of coefficients are bounded.  Translation and matrix
inversion are bounded-degree rational operations on the original
coefficients and the affine point coordinates.  Consequently every
coefficient in \eqref{eq:index-implicit-system} has absolute height
$O(H_{i\ell})$.  Normalizing a nonzero coefficient of each original
equation to one justifies this assertion for coefficient vectors
initially specified only projectively.

Fix an algebraic local frame $e$, and a multi-index $\alpha$ attaining
\eqref{eq:weighted-index} in the factorwise parameters $T_\ell$ just
chosen.  If $a_\alpha$ is its coefficient, then
\[
 a_{\mathrm{fib}}=a_\alpha e(y_i)^{n_i}
       \in (E_i^{n_i})_{y_i}
\]
is a nonzero algebraic fiber element.  In any analytic frame its
coefficient times the value of that frame to the power $n_i$ equals
$a_{\mathrm{fib}}$.  Indeed a change of frame multiplies the
coefficient by its constant term, while all other contributions
involve coefficients of strictly smaller weighted order and vanish.
The multi-index is kept fixed here; only the index itself is asserted
to be independent of a change of parameters.

We will realize the parameters $T_\ell$ on analytic polydisks of radii
$\rho_{i\ell v}$ and choose analytic frames $e_v$ there.  If
\[
 \kappa_v=\sup_z\log\frac{\|e_v(y_i)\|_v}{\|e_v(z)\|_v},
\]
Cauchy's coefficient inequality, or its nonarchimedean Gauss-norm
counterpart, gives
\begin{equation}\label{eq:index-cauchy}
 \log\|a_{\mathrm{fib}}\|_v
 \le \log\|s_i\|_{v,\sup}+n_i\kappa_v
       +\sum_\ell|\alpha_\ell|\log\rho_{i\ell v}^{-1}.
\end{equation}
The estimate is first applied on smaller closed polydisks and then
passed to the indicated radii.  Because the left side uses the same
algebraic fiber element at every place, the product formula says
\[
 \sum_v\nu_v\log\|a_{\mathrm{fib}}\|_v
       =-n_i h_{\bar E_i}(y_i).
\]
It remains to bound the two losses on the right of
\eqref{eq:index-cauchy}: the inverse radii and the metric variation.

\emph{Analytic radii from the conormal equations.}
Let $S_{i\ell v}$ be the maximum of one and the absolute values of
the coefficients of \eqref{eq:index-implicit-system}.  Then
\begin{equation}\label{eq:index-coefficient-sum}
 \sum_v\nu_v\log S_{i\ell v}\ll H_{i\ell}.
\end{equation}
Fix tolerances $0<\gamma_v\le1$, determined by the place of a fixed
field of definition of the ambient data, equal to one outside its
infinite and finitely many finite places.  We may use the radii
\begin{equation}\label{eq:index-radii}
 \rho_{i\ell v}
 =\left(\frac{\gamma_v}{c_vS_{i\ell v}}\right)^4,
 \qquad
 \sum_v\nu_v\log\rho_{i\ell v}^{-1}\le C_\gamma H_{i\ell},
\end{equation}
where $c_v=1$ at finite places and a fixed sufficiently large constant
at infinite places.

Here is the analytic justification, including uniformity over
extensions.  At a finite place, on a smaller closed $T_\ell$-polydisk
of radius $s<\rho_{i\ell v}$, solve
$U=-\beta T-Q(T,U)$ in the Gauss-norm ball $\|U\|\le S_{i\ell v}s$.
The quadratic term has norm at most $S_{i\ell v}^3s^2$, and its
Lipschitz constant in $U$ is at most $S_{i\ell v}^2s<1$.
The map therefore preserves the ball and is a contraction.
At an infinite place use instead the ball
$\|U\|\le C_1S_{i\ell v}s$, with a fixed $C_1$ larger than twice the
maximum number of entries in a row of $\beta_\ell$. This absorbs the
row sums in its linear term. Bounds for the quadratic term and its
Lipschitz constant now acquire only fixed factors depending on $C_1$,
the degrees and the number of monomials. Increasing $c_v$ makes the
quadratic contribution smaller than half the ball radius and makes the
Lipschitz constant less than one.  The solutions for smaller disks agree and give an
analytic parametrization on the open polydisk of radius
$\rho_{i\ell v}$.  All affine coordinate displacements have norm
less than $\gamma_v$, and the Jacobian in $U$ remains invertible.
At the center these equations define the germ of $Z_{i\ell}$:
they define a smooth complete intersection of its dimension containing
that germ.  Every further equation of $Z_{i\ell}$ thus vanishes as
a power series after substitution, and hence on the whole analytic
branch.  The invertible minor shows that this branch remains in the
smooth locus.  Finally \eqref{eq:index-coefficient-sum} gives the
second assertion of \eqref{eq:index-radii}.  Notice that sums of
weights over the places above a fixed place do not change under
extension of $E$.

\emph{Metric variation at the infinite places.}
We next choose the tolerances so that the metric varies little on the
product of these branches.  At infinity, fix a small $\xi>0$.
The bound on affine coordinate displacements gives a uniform bound
on projective distance: for $x=[1:a]$ and $x'=[1:a+u]$, the chordal
distance is at most a fixed constant times $\|u\|$, because both
homogeneous coordinate norms are at least one. Compactness of the fixed
complex abelian varieties and local exponential charts therefore allow
us to make each group difference from the center have an analytic
logarithm of norm at most $\xi$.  Its image
under $F_i$ has a logarithmic lift of norm
$O(\xi\sum_\ell b_{i\ell})=O(\xi M_i)$.  The same coefficient
bound holds at every infinite embedding: all norms on the fixed
finite-dimensional real homomorphism spaces are comparable.

On the target universal cover, centered at a lift of $F_i(y_i)$,
the canonical curvature has a fixed quadratic potential $Q$.
There is a holomorphic frame $e_v$ whose negative log norm differs
from its value at the origin by $Q$: subtracting this potential
leaves a flat metric, with a unitary holomorphic frame on the simply
connected cover.  Pull this frame back along the logarithmic lift
of our polydisk.  We obtain
\begin{equation}\label{eq:index-infinite-metric}
 \sup_z\log\frac{\|e_v(y_i)\|_v}{\|e_v(z)\|_v}
 \le C\xi^2M_i^2.
\end{equation}
The constant is independent of the center.  The image of the polydisk
need not lie in an injectivity neighborhood on the target: the
construction takes place on its universal cover.

\emph{Metric variation at the finite places.}
At a finite place we need a different estimate, independent of all
integer coefficients of $F_i$.  Outside a fixed finite set of places,
the varieties, group laws, polarization data and basis homomorphisms
extend over good models.  Differences of reduction zero remain of
reduction zero under every integral linear combination of the basis
homomorphisms.  The choice $\gamma_v=1$ therefore suffices there.

At one of the remaining places, work over the completed algebraic
closure.  Arbitrarily small analytic balls at the identity are
subgroups.  Indeed, in fixed smooth group parameters addition is
$u+v$ plus terms of degree at least two, and inversion is $-u$ plus
terms of degree at least two; on a sufficiently small ball the
nonarchimedean inequality makes both operations preserve that ball.
In particular the ball is stable under multiplication by every
integer.  Choose a target subgroup ball small enough that translation
by any of its points preserves the reduction class in the fixed
projective embedding.  Such a choice is uniform in the point being
translated.  Here is a uniform projective estimate for this assertion. For normalized
homogeneous coordinates, write
\[
 d_v(x,x')=\frac{\max_{a,b}|x_ax'_b-x_bx'_a|_v}
                  {\max_a|x_a|_v\max_a|x'_a|_v}.
\]
In a common affine chart, coordinate displacements of norm less than
$\gamma_v\le1$ imply $d_v(x,x')<\gamma_v$, even when the affine
coordinates of the center are unbounded. A fixed projective morphism $f$
has finitely many homogeneous polynomial presentations which together
have no base point on its source. A homogeneous Nullstellensatz
identity gives a fixed $\varepsilon_v>0$ such that at every normalized
source point one presentation has value of norm at least $\varepsilon_v$. Polynomial
difference estimates give a fixed $C_v'$ such that, when
$d_v(x,x')$ is below a fixed threshold, the same presentation remains
nonzero and
\[
 d_v(f(x),f(x'))\le C_v' d_v(x,x').
\]
To apply the polynomial estimate, rescale nearby normalized coordinate
vectors so that their coordinate differences are bounded by their
projective distance; this is possible in a coordinate which is a unit
at both points. All constants are in the fixed coefficient field and
remain valid over its completed algebraic closure.
Apply the estimate to the difference morphism and to translation
$A\times A\to A$, comparing $(a,b)$ with $(a,0)$. Thus the source
coordinate tolerances uniformly control identity neighborhoods, and
translation by a sufficiently small identity neighborhood preserves
projective reduction for every center.

The finitely many basis homomorphisms carry sufficiently small source
identity neighborhoods into this target subgroup ball.  Every integral
linear combination does so too.  We can therefore choose the input
tolerances $\gamma_v$ independently of $i$ so that the target remains
in the reduction class of $F_i(y_i)$ throughout the product branch.
A homogeneous coordinate having maximal absolute value at this center
defines a nonvanishing frame with constant projective metric norm on
the branch.  The comparison between this metric and the canonical
metric gives
\begin{equation}\label{eq:index-finite-metric}
 \sup_z\log\frac{\|e_v(y_i)\|_v}{\|e_v(z)\|_v}\le C_v,
 \qquad \sum_{v\nmid\infty}\nu_v C_v\le C_0,
\end{equation}
where $C_v=0$ outside the fixed bad places.  These constants remain
unchanged upon extension of the point field.  This is the step that
prevents a loss depending on the sizes of the integers in $F_i$.

\emph{Summing over the places.}
The constructed frames satisfy $\kappa_v\le C\xi^2M_i^2$ at infinity
and $\kappa_v\le C_v$ at finite places.  The section norm bound applies
on all the branches because they lie in the resolution's isomorphic
smooth locus.  Sum \eqref{eq:index-cauchy}, using the product formula,
the small-section bound \eqref{eq:index-small-section}, the radius
estimate \eqref{eq:index-radii}, and the two metric estimates
\eqref{eq:index-infinite-metric} and \eqref{eq:index-finite-metric}.  We
obtain
\[
 C_\gamma\sum_\ell|\alpha_\ell|H_{i\ell}
 \ge n_i\bigl((c-C\xi^2)M_i^2-C_0-o(M_i^2)\bigr).
\]
Choose $\xi$ so that $C\xi^2<c/2$, then let $i$ tend to infinity.
The comparison $b_{i\ell}^2H_{i\ell}\asymp M_i^2$ now yields a
fixed positive lower bound for \eqref{eq:weighted-index}.
\end{proof}

For the products chosen in Proposition~\ref{prop:minimal-products},
the first hypothesis is precisely its smoothness and equation
conclusion.  The homomorphisms and weights of
Section~\ref{sec:arithmetic} satisfy the second hypothesis;
\eqref{eq:small-marked-height} and
Lemma~\ref{lemma:small-sections} give the third.  We therefore have
a fixed positive index at $y_i$.  The next section converts it into
a hypersurface of bounded degree over a field of bounded degree,
which is the contradiction required by minimality.
